%% kapa: exam=1 round=29 year=2018 period=1 form=A subject=civil_law subject_ko=민법 seq=23 number=23
%% source: Q-Net 감정평가사 기출문제 (cst00309) artlSeq=5202808 file=extracted/1008/5202808/2018년 제29회 감정평가사 1차시험 기출문제/2018년 제29회 감정평가사 1차시험 기출문제/1교시 A형.pdf
%% conversion: kordoc 4.24.51 -> tools/gen1.py (auto); review=proofread-sonnet; flags=none
\begin{question}{23}
전세권에 관한 설명으로 옳은 것을 모두 고른 것은? (다툼이 있으면 판례에 따름)
\begin{framedtext}
ㄱ. 전세권자는 전세권이 설정된 부동산 전부에 대하여 후순위 권리자나 그 밖의 일반채권자보다 전세금의 우선변제를 받을 권리가 있다. \par
ㄴ. 전세권은 용익물권적 성격과 담보물권적 성격을 겸비하고 있다. \par
ㄷ. 타인의 토지에 있는 건물에 전세권을 설정한 경우 전세권의 효력은 그 건물의 소유를 목적으로 한 지상권에는 미치지 않는다. \par
\end{framedtext}
\begin{choices}
  \item ㄱ
  \item ㄷ
  \item ㄱ, ㄴ
  \item ㄴ, ㄷ
  \item ㄱ, ㄴ, ㄷ
\end{choices}
\end{question}
\answer{3}
